\chapter*{How to read this guide}

This Blueprint is a guided reading map of the main proof route.  Its nodes group declarations
by mathematical purpose, so it is intentionally smaller than the source.  The
\href{dependency-graph/index.html}{scoped interactive declaration graph} is generated from
the compiled Lean environment; generated declaration paths are collapsed and LU modules are
shown as boundary nodes.  Its
\href{dependency-graph/graph.json}{raw JSON} records the same graph for other tools.  The
interactive view initially hides private declarations; it can reveal them, reverse edge
direction, filter edge kinds, and focus on either side of a selected declaration.

\chapter{Echelon predicates}

\begin{definition}[Echelon and reduced-echelon predicates]
\label{def:echelon-predicates}
\lean{Matrix.RowIsZero, Matrix.IsPivot, Matrix.IsEchelonForm,
  Matrix.IsReducedEchelonForm}
\leanok
These predicates are the semantic target of the executable algorithm.  They separate the
meaning of a valid output from the control flow that computes it, so later correctness results
can expose stable mathematical statements.
\end{definition}

\begin{theorem}[Zeros below a pivot]
\label{thm:pivot-column-zero-below}
\lean{Matrix.IsEchelonForm.pivot_column_zero_below}
\leanok
\uses{def:echelon-predicates}
An echelon witness forces every entry below a pivot to vanish.  This is the main local fact that
downstream proofs can consume without unfolding the inductive echelon predicate.
\end{theorem}

\chapter{Row-operation data and executable algorithm}

\begin{definition}[Row-operation data]
\label{def:row-operation-data}
\lean{RowOp, swapRow, factor, replace, checkPivot, Matrix.RowReductionResult}
\leanok
The primitive row operations, pivot search, and result log form the executable vocabulary.
Keeping the log beside the resulting matrix makes later replay and preservation theorems
possible without changing the public matrix-returning wrappers.
\end{definition}

\begin{definition}[Single-column elimination loop]
\label{def:elimination-loop}
\lean{GaussianEliminationInternal.eliminateColLoopAux,
  GaussianEliminationInternal.eliminateColLoop,
  GaussianEliminationInternal.eliminateColCore}
\leanok
\uses{def:row-operation-data}
This loop chooses and normalizes a pivot, eliminates the selected column, and records every row
operation.  Isolating one column gives the correctness proof a small transition to analyze.
\end{definition}

\begin{definition}[REF and RREF algorithms]
\label{def:row-reduction-algorithm}
\lean{GaussianEliminationInternal.rowReductionAux,
  GaussianEliminationInternal.rawRowEchelonForm,
  GaussianEliminationInternal.rawReducedRowEchelonForm,
  Matrix.rowEchelonForm, Matrix.reducedRowEchelonForm}
\leanok
\uses{def:row-operation-data, def:elimination-loop}
The outer recursion advances through rows and columns, while the public wrappers discard the
operation log.  The effect is a computable REF or RREF API whose proof-oriented internals still
retain enough evidence for correctness and reconstruction.
\end{definition}

\chapter{Row-equivalence semantics}

\begin{definition}[Row equivalence]
\label{def:row-equivalence}
\lean{Matrix.RowEquivalent, Matrix.RowEquivalent.refl, Matrix.RowEquivalent.symm,
  Matrix.RowEquivalent.trans, Matrix.rowEquivalentSetoid}
\leanok
Row equivalence packages the existence of an invertible left multiplier and provides the
equivalence laws.  It is the semantic bridge from low-level row updates to statements about
the linear system represented by a matrix.
\end{definition}

\begin{theorem}[Elementary row-operation semantics]
\label{thm:elementary-row-semantics}
\lean{Matrix.applyRowOp, Matrix.elementaryMatrixOfRowOp,
  Matrix.elementaryMatrixOfRowOp_mul_eq_applyRowOp,
  Matrix.swap_matrix_eq_elem_mul_matrix, Matrix.factor_matrix_eq_elem_mul_matrix,
  Matrix.replace_matrix_eq_elem_mul_matrix, Matrix.swap_inv, Matrix.factor_inv,
  Matrix.replace_inv,
  Matrix.rowEquivalent_swapRow, Matrix.rowEquivalent_factor,
  Matrix.rowEquivalent_replace}
\leanok
\uses{def:row-operation-data, def:row-equivalence}
Each executable row operation agrees with left multiplication by its elementary matrix, and
each operation preserves row equivalence.  This turns an algorithmic step into algebra that can
be composed across the whole reduction.
\end{theorem}

\begin{definition}[An echelon representative]
\label{def:echelon-representative}
\lean{Matrix.IsEchelonFormOf, Matrix.IsEchelonFormOf.rowEquivalent,
  Matrix.IsEchelonFormOf.echelon, Matrix.isEchelonFormOf_mk}
\leanok
\uses{def:echelon-predicates, def:row-equivalence}
This bundled relation says that a matrix is both echelon and row-equivalent to the source.
Bundling the two obligations gives later clients one reusable certificate instead of parallel
hypotheses.
\end{definition}

\begin{theorem}[Column-operation semantics]
\label{thm:column-operation-semantics}
\lean{ColumnElementary.swapCol, ColumnElementary.factorCol,
  ColumnElementary.replaceCol, ColumnElementary.swapCol_eq_mul_elem,
  ColumnElementary.factorCol_eq_mul_elem, ColumnElementary.replaceCol_eq_mul_elem}
\leanok
\uses{def:row-operation-data}
Transpose-derived column operations are characterized as right multiplication by elementary
matrices.  They are included because LU construction consumes the row-operation log from the
opposite side.
\end{theorem}

\chapter{Pivot and state-machine correctness}

\begin{theorem}[Pivot search]
\label{thm:pivot-search}
\lean{Matrix.checkPivot_none_zero, Matrix.checkPivot_some_row_ge,
  Matrix.checkPivot_some_minimal, Matrix.checkPivot_some_nonzero}
\leanok
\uses{def:row-operation-data}
The search result is either evidence that the candidate column is zero below the current row or
a minimal nonzero pivot at a valid row.  These facts justify both branches of the executable
state transition.
\end{theorem}

\begin{theorem}[Column-elimination invariants]
\label{thm:column-elimination-invariants}
\lean{Matrix.eliminateCol_matrix_irrel, Matrix.eliminateCol_go_preserves_col,
  Matrix.eliminateCol_pivotRow, Matrix.eliminateCol_pivotCol_zero,
  Matrix.eliminateCol_below_pivotCol_zero}
\leanok
\uses{def:elimination-loop}
One elimination transition preserves earlier columns, establishes the pivot row, and clears the
new pivot column below it.  These are exactly the local invariants needed to extend an echelon
prefix by one step.
\end{theorem}

\begin{definition}[Echelon machine state]
\label{def:echelon-state}
\lean{Matrix.EchelonStateCore, Matrix.ReducedStateExtras, Matrix.EchelonState,
  Matrix.echelon_of_state, Matrix.reduced_of_state}
\leanok
\uses{def:echelon-predicates}
The state records which prefix has already been justified and adds the extra conditions needed
for reduced echelon form.  Its projections turn accumulated loop invariants back into the
public semantic predicates.
\end{definition}

\begin{theorem}[Algorithmic echelon correctness]
\label{thm:echelon-correctness}
\lean{Matrix.rowEchelonForm_isEchelonForm,
  Matrix.reducedRowEchelonForm_isEchelonForm,
  Matrix.reducedRowEchelonForm_isReducedEchelonForm}
\leanok
\uses{def:row-reduction-algorithm, thm:column-elimination-invariants,
  def:echelon-state, thm:pivot-search}
The terminal machine state proves that the executable wrappers satisfy their advertised REF and
RREF predicates.  Users can rely on these theorems without inspecting recursion or pivot cases.
\end{theorem}

\chapter{Log and solution preservation}

\begin{theorem}[Solution preservation for primitive steps]
\label{thm:primitive-solution-preservation}
\lean{Matrix.mul_by_inv, Matrix.swap_proof, Matrix.factor_proof,
  Matrix.replace_proof}
\leanok
\uses{def:row-operation-data, thm:elementary-row-semantics}
Invertibility of elementary operations shows that multiplying a vector to zero is unchanged by
each primitive step.  This is the local preservation law used to lift solution equivalence
through a recorded reduction.
\end{theorem}

\begin{theorem}[Solution preservation for complete reductions]
\label{thm:reduction-solution-preservation}
\lean{Matrix.eliminate_proof, Matrix.ref_proof, Matrix.rref_proof}
\leanok
\uses{def:row-reduction-algorithm, thm:primitive-solution-preservation}
The local preservation laws are composed over the elimination recursion.  The result is that
REF and RREF compute different matrix representatives without changing the homogeneous solution
set.
\end{theorem}

\begin{theorem}[Operation-log replay]
\label{thm:operation-log-replay}
\lean{Matrix.elim_col_adds_steps, Matrix.row_reduction_adds_steps,
  Matrix.rowEchelonForm_steps, Matrix.reducedRowEchelonForm_steps}
\leanok
\uses{def:row-reduction-algorithm, thm:elementary-row-semantics}
Replaying the recorded operations with \texttt{applyRowOp} reconstructs the computed matrix.
This connects the audit log to executable output and supplies the bridge consumed by LU
factorization.
\end{theorem}

\chapter{RREF semantic uniqueness}

\begin{definition}[A reduced-echelon representative]
\label{def:reduced-echelon-representative}
\lean{Matrix.IsReducedEchelonFormOf,
  Matrix.IsReducedEchelonFormOf.rowEquivalent,
  Matrix.IsReducedEchelonFormOf.reduced,
  Matrix.isReducedEchelonFormOf_mk}
\leanok
\uses{def:echelon-predicates, def:row-equivalence}
This is the reduced analogue of an echelon representative: it combines row equivalence to the
source with the reduced predicate.  It states uniqueness at the semantic boundary rather than
in terms of one implementation.
\end{definition}

\begin{theorem}[The algorithm produces the semantic representative]
\label{thm:algorithmic-rref-representative}
\lean{Matrix.reducedRowEchelonForm_rowEquivalent,
  Matrix.reducedRowEchelonForm_isReducedEchelonFormOf,
  Matrix.reducedRowEchelonForm_isEchelonFormOf}
\leanok
\uses{def:row-reduction-algorithm, thm:echelon-correctness,
  def:reduced-echelon-representative}
The executable RREF is row-equivalent to its input and satisfies the reduced predicate.  This
connects the algorithmic result to the implementation-independent uniqueness statement.
\end{theorem}

\begin{definition}[Canonical algorithmic RREF]
\label{def:canonical-rref}
\lean{Matrix.IsCanonicalRrefOf,
  Matrix.isCanonicalRrefOf_reducedRowEchelonForm,
  Matrix.IsCanonicalRrefOf.unique, Matrix.IsCanonicalRrefOf.reduced,
  Matrix.reducedRowEchelonForm_isCanonicalRrefOf}
\leanok
\uses{def:row-reduction-algorithm, thm:echelon-correctness}
The canonical relation identifies a matrix with the implementation's RREF output.  Its immediate
uniqueness is useful for executable equality checks, while the reduced theorem connects it back
to the semantic predicate.
\end{definition}

\begin{theorem}[Semantic uniqueness of RREF]
\label{thm:semantic-rref-uniqueness}
\lean{Matrix.IsReducedEchelonFormOf.unique,
  Matrix.IsReducedEchelonFormOf.canonical,
  Matrix.RrefUniquenessSemanticGoal,
  Matrix.rrefUniquenessSemanticGoal_holds}
\leanok
\uses{def:reduced-echelon-representative,
  thm:algorithmic-rref-representative, thm:pivot-column-zero-below}
Any two reduced-echelon representatives of the same source matrix are equal.  This closes the
proof route by upgrading a concrete algorithmic normal form into a semantic uniqueness theorem.
\end{theorem}

\chapter{LU downstream consumers}

\begin{definition}[LU factorization boundary]
\label{def:lu-factorization-boundary}
\lean{Matrix.luFactorization}
\leanok
\uses{def:row-reduction-algorithm, thm:column-operation-semantics}
The LU implementation consumes REF and its operation log to build permutation and lower factors.
This node is a boundary summary: LU declarations are not expanded in the complete Gaussian
graph.
\end{definition}

\begin{theorem}[LU correctness boundary]
\label{thm:lu-correctness-boundary}
\lean{Matrix.luFactorization_reconstruct,
  Matrix.luFactorization_upper_isEchelonForm,
  Matrix.luFactorization_lower_isUnitLowerTriangular}
\leanok
\uses{def:lu-factorization-boundary, thm:echelon-correctness,
  thm:operation-log-replay}
The downstream correctness layer reconstructs the input and certifies structural properties of
the factors.  It demonstrates the effect of the Gaussian theory while deliberately treating LU
internals as outside this visualization's scope.
\end{theorem}
